\section*{Chapitre \ref{ch:mx:smart-order-routing} --- Routage intelligent d'ordres}

\begin{solution}{exo:mx:smart-order-routing:1}
La plateforme la plus éloignée est à 360~microsecondes~: le fragment envoyé en A attend $360-50=310$ microsecondes, celui envoyé en C $360-200=160$.
\end{solution}

\begin{solution}{exo:mx:smart-order-routing:2}
L'\omterm{def:mx:smart-order-routing:spray}{ordre dispersé} s'exécute encore en premier en A (50~microsecondes). Le fournisseur en est informé 95~microsecondes plus tard, à 145~; ses annulations atteignent C à 150, B à 195, D à
195 et E à 245. L'\omterm{def:mx:smart-order-routing:spray}{ordre dispersé} atteint B à 120, avant l'annulation, et C, D et E à 200, 280 et 360, après elle~: il trouve les cotations rapides en A et B,
1\,600 des 2\,500 actions (64~\%).
\end{solution}

\begin{solution}{exo:mx:smart-order-routing:3}
Au cours vendeur, 1\,300 actions ne paient que des frais~: $500\times0.30+200\times(0.30+0.10+0.30-0.10)=270$ cents. Le balayage paie un cent de plus sur 1\,200 actions
(1\,200 cents) et des frais de $1\,000\times0.30+200\times0.30=360$ cents. Au-dessus de 100.01~: $(270+1\,200+360)/2\,500=0.73$
cents par action.
\end{solution}

\begin{solution}{exo:mx:smart-order-routing:4}
La course est perdue parce que l'exécution en A arrive toujours en premier et que les annulations sont plus rapides que les fragments suivants du routeur. Des latences aléatoires font parfois
arriver un fragment éloigné avant que l'exécution en A ait été signalée, et trouver alors la cotation rapide~; ces latences aident l'\omterm{def:mx:smart-order-routing:spray}{ordre dispersé} exactement comme elles nuisent au balayage synchronisé.
\end{solution}

\begin{solution}{exo:mx:smart-order-routing:5}
Dans la plupart des scénarios, chaque plateforme n'exécute qu'une partie de son ordre, si bien qu'afficher plus que nécessaire augmente le nombre espéré d'actions exécutées vers $X$. La pénalité de
sur-exécution (1.2 cents par action, plus que les 0.7 environ que rapporte chaque exécution passive) rend coûteux les scénarios où la plupart des plateformes exécutent, ce qui borne la
surréservation (1\,778 actions affichées).
\end{solution}

\begin{solution}{exo:mx:smart-order-routing:6}
Il affiche tout sur la plateforme offrant un rabais de 0.25 cent, qui a la file la plus longue (2\,500 actions)~: seulement 22.0~\% s'exécutent, et la pénalité pour le
reste porte son coût à 0.77 cent par action, presque autant qu'un \omterm{def:mx:the-limit-order-book:orders}{ordre au marché}.
\end{solution}

\begin{solution}{exo:mx:smart-order-routing:7}
Afficher paie~: chaque exécution passive rapporte plus que la pénalité de sur-exécution de 0.5 cent. L'optimum affiche environ 4\,030 actions (1\,000 en A, B et D, 751 en C,
277 au dark pool) et «~coûte~» $-0.24$ cent par action hors échantillon, un gain qui n'existe que parce que la pénalité sous-estime le risque d'acheter
trop.
\end{solution}

\begin{solution}{exo:mx:smart-order-routing:8}
Les 500 actions de E étaient là quand le routeur a regardé~; 300 d'entre elles appartenaient à un fournisseur qui a annulé après que l'exécution du routeur en A l'a atteint.
Le routeur, par son propre \omterm{def:mx:smart-order-routing:spray}{ordre dispersé}, est à l'origine de la perte. Synchronisé, E s'exécute à 100~\% et, avec le rabais du preneur de liquidité, devient la plateforme la moins chère.
\end{solution}

\begin{solution}{pb:mx:smart-order-routing:1}
\textbf{1.} Il découpe un ordre entre plusieurs plateformes et choisit la plateforme, le prix, la taille, le type et le moment de chaque fragment.
\textbf{2.} De A à E, à 50, 120, 200, 280 et 360 microsecondes~; frais preneur de 0.30, 0.30, 0.10, 0.30 et $-0.10$ cent~; chacune affiche 200 actions lentes et 300
actions rapides à 100.01 et 1\,000 à 100.02~; le fournisseur rapide se trouve en A, avec des liaisons de 45 à 195 microsecondes, et annule partout après une exécution.
\textbf{3.} Voir les définitions.
\textbf{4.} Le carnet consolidé était plus profond après l'arrivée du nouvel entrant, et il fournissait moins de liquidité là où ses meilleures cotations étaient traversées.
\textbf{5.} L'\omterm{def:mx:smart-order-routing:spray}{ordre dispersé} arrive à 50, 120, 200, 280, 360~; les annulations à 100, 150, 200 et 250 de B à E.
\textbf{6.} 52~\%, et 0.73 cent par action au-dessus du cours vendeur avec le balayage et les frais.
\textbf{7.} Chaque fragment attend la latence la plus lointaine moins la sienne~: 310, 240, 160, 80 et 0 microsecondes.
\textbf{8.} \emph{Résultat nommé}~: un \omterm{def:mx:smart-order-routing:spray}{ordre dispersé} s'exécute à 52~\% des 2\,500 actions affichées et coûte 0.73 cent par action, net des frais~; synchronisé, il s'exécute
à 100~\% pour un coût de 0.18 cent, soit le frais moyen payé par le preneur.
\textbf{9.} Jusqu'à 10~\% de gigue, rien ne change~; à 20~\%, le balayage synchronisé s'exécute à 99.3~\% (la totalité du niveau dans 94~\% des essais) et à 40~\%, il s'exécute à 90.3~\%~; pour l'\omterm{def:mx:smart-order-routing:spray}{ordre
dispersé}, il monte à 68.6~\%.
\textbf{10.} A 100~\%, B 44~\%, C 41~\%, D 40~\%, E 40~\%.
\textbf{11.} \omterm{def:mx:smart-order-routing:spray}{Ordre dispersé}~: A, E, C, B, D (de 0.30 à 0.72 cent)~; synchronisé~: E, C, puis A, B, D ($-0.10$ à 0.30)~: un classement mesure le
routage propre au routeur.
\textbf{12.} Par frais (E, C)~: 99.8~\% en \omterm{def:mx:smart-order-routing:spray}{ordre dispersé}, 100~\% en synchronisé, coût voisin de zéro~; par distance (A, B)~: 71.2~\% en \omterm{def:mx:smart-order-routing:spray}{ordre dispersé} à 0.59 cent, 100~\% en synchronisé à
0.30.
\textbf{13.} Sa première exécution venait de plateformes lointaines, à partir desquelles les annulations du fournisseur rapide ne pouvaient pas dépasser l'autre fragment.
\textbf{14.} $(h+f)M-\sum_k(h+r_k)\min(L_k,(\xi_k-Q_k)^+)+\lambda_u(X-A)^++\lambda_o(A-X)^+$.
\textbf{15.} 0.80 (100~\%), 0.77 (22.0~\%), 0.74 (38.7~\%), 0.20 (61.3~\%) et 0.03 cents (76.4~\%).
\textbf{16.} 0.03 cent par action (0.06 sans cela).
\textbf{17.} Les \omterm{def:mx:the-limit-order-book:orders}{ordres à cours limité} envoyés aux plateformes versant les plus gros rabais ont été moins bien exécutés~: router pour les rabais ne maximise pas la qualité d'exécution.
\textbf{18.} Arrivées synchronisées, avec \omterm{def:mx:smart-order-routing:spray}{évanouissement de la cotation} (quote fade) après une première exécution, \omterm{def:mx:dark-pools-and-blocks:minqty}{quantités minimales d'exécution} dans les dark pools, \omterm{def:mx:dark-pools-and-blocks:ping}{ordres sondes} (pinging),
tailles et moments des \omterm{def:mx:the-almgren-chriss-framework:parent}{ordres enfants} randomisés (détection de motifs).
\textbf{19.} Synchroniser les balayages, classer les plateformes sur les exécutions mesurées et non sur les affichages, et répartir les ordres passifs selon les exécutions espérées et non selon les rabais.
\textbf{20.} Autour de cotations qui ne seront pas là et de files qui ne s'écouleront pas.
\end{solution}

\begin{solution}{iq:mx:smart-order-routing:1}
\omterm{def:mx:smart-order-routing:spray}{Évanouissement de la cotation} (quote fade)~: les fragments sont arrivés les uns après les autres, les premières exécutions ont alerté les fournisseurs rapides, qui ont annulé ailleurs. Vérifier les heures d'arrivée par plateforme et
les annulations qui les précèdent immédiatement~; synchroniser.
\iqlookfor{Écart des heures d'arrivée~; réaction du fournisseur.}
\end{solution}

\begin{solution}{iq:mx:smart-order-routing:2}
Mesurer la latence à sens unique par plateforme (à partir des accusés de réception, en continu), retenir chaque fragment de l'écart avec la plus lente, et ajouter une marge pour la gigue.
Elle échoue lorsque les latences varient plus que le temps de réaction du fournisseur le plus rapide, ou lorsque la passerelle d'une plateforme met les messages en file.
\iqlookfor{Estimation de la latence~; gigue comme limite.}
\end{solution}

\begin{solution}{iq:mx:smart-order-routing:3}
À partir des ordres propres du routeur~: exécutions selon la durée de repos, conditionnelles à la position dans la file et à la taille~; la difficulté tient à la censure (ordres annulés),
la sélection (le routeur n'affiche qu'aux endroits où il espère des exécutions) et la dépendance à son propre comportement.
\iqlookfor{Censure et sélection~; conditionnement sur la file.}
\end{solution}

\begin{solution}{iq:mx:smart-order-routing:4}
Choisir un \omterm{def:mx:the-limit-order-book:orders}{ordre au marché} $M$ et des \omterm{def:mx:the-limit-order-book:orders}{ordres à cours limité} $L_k$~; chaque \omterm{def:mx:the-limit-order-book:orders}{ordre à cours limité} s'exécute sur ce qui sort de la tête de sa file au-delà de la file qui le précède, dans la limite de $L_k$~; minimiser le
coût espéré net des frais, plus les pénalités de sous-exécution et de sur-exécution (Cont et Kukanov), résolu sur des scénarios ou par approximation stochastique.
\iqlookfor{Exécution comme $\min(L,(\xi-Q)^+)$~; pénalités~; convexité.}
\end{solution}

\begin{solution}{iq:mx:smart-order-routing:5}
Répondre en exposant la logique de routage (exécutions et coût espérés, et non rabais) et la montrer~: taux d'exécution, délais d'exécution et markouts par plateforme pour les ordres propres du
client, face aux rabais, comme le ferait l'analyse de Battalio, Corwin et Jennings.
\iqlookfor{Conflit d'intérêts~; preuves par plateforme.}
\end{solution}

\begin{solution}{iq:mx:smart-order-routing:6}
Par plateforme, groupe de symboles et type d'ordre~: actions envoyées, actions exécutées, délai d'exécution, frais, latences et markouts~; pondérés selon une demi-vie, avec une loi a priori pour
les nouvelles plateformes~; le routeur classe les plateformes selon le coût espéré, y compris les non-exécutions, et explore un peu pour garder les estimations à jour.
\iqlookfor{Champs~; vieillissement~; exploration.}
\end{solution}
